%% GENERATED by tools/from_docs.py from stepss-docs. Do not edit by hand:
%% edit the page named below in stepss-docs and regenerate.
%% source: stepss-docs/src/content/docs/test-systems/index.md

\chapter{Test systems}\label{doc:ts-index}

Four benchmark networks ship as ready-to-run RAMSES data sets, each in its own
repository. They are ordered here from smallest to largest, which is also the
order in which they are worth meeting.

{\def\LTcaptype{none} % do not increment counter
\begin{small}\begin{longtable}[]{@{}
  >{\raggedright\arraybackslash}p{(\linewidth - 8\tabcolsep) * \real{0.1739}}
  >{\raggedleft\arraybackslash}p{(\linewidth - 8\tabcolsep) * \real{0.1522}}
  >{\raggedleft\arraybackslash}p{(\linewidth - 8\tabcolsep) * \real{0.2174}}
  >{\raggedright\arraybackslash}p{(\linewidth - 8\tabcolsep) * \real{0.2391}}
  >{\raggedright\arraybackslash}p{(\linewidth - 8\tabcolsep) * \real{0.2174}}@{}}
\toprule\noalign{}
\begin{minipage}[b]{\linewidth}\raggedright
System
\end{minipage} & \begin{minipage}[b]{\linewidth}\raggedleft
Buses
\end{minipage} & \begin{minipage}[b]{\linewidth}\raggedleft
Machines
\end{minipage} & \begin{minipage}[b]{\linewidth}\raggedright
Frequency
\end{minipage} & \begin{minipage}[b]{\linewidth}\raggedright
Best for
\end{minipage} \\
\midrule\noalign{}
\endhead
\bottomrule\noalign{}
\endlastfoot
\hyperref[doc:ts-5bus]{5-Bus} & 5 & 1 & 50 Hz & Learning the tools \\
\hyperref[doc:ts-kundur]{Kundur Two-Area} & 11 & 4 & 60 Hz & Inter-area oscillations, PSS tuning \\
\hyperref[doc:ts-nordic]{Nordic} & 74 & 20 & 50 Hz & Long-term voltage stability \\
\hyperref[doc:ts-gb]{GB Network} & 87 & 37 & 50 Hz & HVDC, wide-area monitoring \\
\end{longtable}\end{small}
}

\section{Which One to Start With}\label{ts-index:which-one-to-start-with}

\begin{itemize}
\tightlist
\item
  \textbf{New to STEPSS}: the \hyperref[doc:ts-5bus]{5-bus system} is small enough to
  trace every computation by hand while still exercising the whole workflow, from
  power flow through disturbance to trajectory extraction.
\item
  \textbf{Small-signal and damping studies}: the
  \hyperref[doc:ts-kundur]{Kundur two-area system} ships two dynamic data variants,
  with and without PSS, so the inter-area mode can be compared directly. Pair it
  with \hyperref[doc:eigenanalysis]{Eigenanalysis}.
\item
  \textbf{Voltage stability and long-term dynamics}: the
  \hyperref[doc:ts-nordic]{Nordic system} is the IEEE PES-TR19 benchmark, with
  several operating points and a Jupyter tutorial that walks through a voltage
  collapse.
\item
  \textbf{Converter-dominated and HVDC studies}: the
  \hyperref[doc:ts-gb]{GB network} models its interconnections as two-port
  converters and carries ready-made disturbance templates.
\end{itemize}

\section{Running Any of Them}\label{ts-index:running-any-of-them}

All four follow the same pattern: clone the repository, then point a stepss
\texttt{cfg} at its data files.

\begin{Verbatim}
import stepss

case = stepss.cfg()
case.addData('dyn.dat')            # dynamic data
case.addData('lf.dat')             # power-flow solution
case.addData('solveroptions.dat')  # solver settings
case.addDst('disturb.dst')         # disturbance scenario
case.addObs('obs.dat')             # observables to record

sim = stepss.sim()
sim.execSim(case)
\end{Verbatim}

The exact file names differ per system; each page lists them. See
\hyperref[doc:py-examples]{Python API Examples} for complete worked scripts, and
\hyperref[doc:py-install]{Installation} if stepss is not set up yet.

\section{Next Steps}\label{ts-index:next-steps}

\begin{itemize}
\tightlist
\item
  \hyperref[doc:py-examples]{Python API Examples}, full simulation workflows
\item
  \hyperref[ch-disturb]{Disturbances}, write your own scenarios
\item
  \hyperref[doc:model-index]{Model Reference}, the models these systems use
\end{itemize}
